% Ce fichier est prévu pour être inclus dans un document LaTeX à deux colonnes.
% Packages requis dans le préambule :
% \usepackage{amsmath,amssymb,array,booktabs,tabularx}

\section{Protocole d'évaluation technique de GEAR}
\label{sec:protocole-evaluation-technique}

L'évaluation technique de GEAR poursuit deux objectifs complémentaires.
Le premier consiste à déterminer dans quelle mesure son métamodèle représente
les capacités des frameworks étudiés et dans quelle mesure ses connecteurs
préservent ces capacités lors de la génération. Le second consiste à vérifier
que les artefacts générés sont valides, exécutables et conformes aux workflows
sources. Ces objectifs sont formulés par les questions de recherche suivantes :

\begin{quote}
\textbf{RQ1 :} Dans quelle mesure GEAR permet-il de représenter et de préserver
les principales capacités de conception et d'orchestration de frameworks
hétérogènes pour systèmes multi-agents basés sur des \gls{llm} ?
\end{quote}

\begin{quote}
\textbf{RQ1.1 :} Dans quelle mesure le métamodèle de GEAR permet-il de
représenter les capacités proposées par les frameworks sélectionnés ?
\end{quote}

\begin{quote}
\textbf{RQ1.2 :} Dans quelle mesure les spécifications GEAR sont-elles
sémantiquement préservées lors de la génération d'implémentations propres aux
différents frameworks ?
\end{quote}

\begin{quote}
\textbf{RQ2 :} Dans quelle mesure les implémentations générées par GEAR
sont-elles exécutables et préservent-elles la sémantique des workflows
spécifiés ?
\end{quote}

\begin{quote}
\textbf{RQ2.1 :} Dans quelle mesure GEAR génère-t-il des artefacts
syntaxiquement valides, importables et exécutables ?
\end{quote}

\begin{quote}
\textbf{RQ2.2 :} Dans quelle mesure les traces des artefacts générés
respectent-elles les contraintes des spécifications GEAR ?
\end{quote}

La suite présente d'abord la sélection des frameworks, puis le protocole de
couverture fonctionnelle associé à RQ1 et, enfin, l'évaluation de la correction
et de la fidélité associée à RQ2.

\section{Sélection des frameworks étudiés}
\label{sec:selection-frameworks}

Les frameworks sélectionnés pour cette étude sont CrewAI, Google ADK,
LangGraph, OpenAI Agents SDK, Microsoft Agent Framework, Strands
Agents, PydanticAI, AutoGen, Semantic Kernel et Haystack.

Nous avons adopté une stratégie d'échantillonnage raisonné visant à
maximiser la diversité des frameworks étudiés. Notre objectif n'était
pas d'établir un classement exhaustif des solutions existantes, mais
de couvrir différentes approches de conception et d'orchestration des
\gls{mas} basés sur des \gls{llm}. Les frameworks retenus représentent
notamment les systèmes organisés autour de rôles et de tâches, les
équipes conversationnelles, les workflows fondés sur des graphes, les
définitions typées d'agents, les architectures en pipeline et les SDK
d'agents proposés par différents fournisseurs.

Un framework a été inclus lorsqu'il répondait aux critères suivants :

\begin{itemize}
    \item il était publiquement accessible et suffisamment documenté ;
    \item il proposait une API Python ;
    \item il permettait de définir au moins un agent basé sur un
    \gls{llm} ;
    \item il permettait de composer plusieurs agents ou plusieurs
    tâches ;
    \item il disposait d'une version identifiable et installable ;
    \item il était activement maintenu au moment de la sélection ;
    \item il représentait un paradigme pertinent de conception ou
    d'orchestration des systèmes multi-agents.
\end{itemize}

À l'inverse, nous avons exclu les bibliothèques limitées aux appels
directs à des \gls{llm}, les plateformes fermées ne proposant pas
d'API accessible, les projets inactifs, les frameworks exclusivement
destinés aux agents d'apprentissage par renforcement et les outils
d'observabilité ne permettant pas de construire des systèmes
multi-agents. Dans la mesure où la version actuelle de GEAR génère des
implémentations en Python, les frameworks exclusivement destinés à
d'autres langages de programmation ont également été exclus.

Le Tableau~\ref{tab:frameworks-selectionnes} présente le principal
paradigme représenté par chacun des frameworks sélectionnés.

\begin{table}[t]
    \centering
    \caption{Frameworks sélectionnés et principaux paradigmes représentés.}
    \label{tab:frameworks-selectionnes}
    \small
    \begin{tabularx}{\columnwidth}{@{}p{0.28\columnwidth}>{\raggedright\arraybackslash}X@{}}
        \toprule
        \textbf{Framework} & \textbf{Principal paradigme représenté} \\
        \midrule
        CrewAI & Organisation des agents autour de rôles et de tâches \\
        Google ADK & Agents hiérarchiques et agents de workflow \\
        LangGraph & Orchestration par graphe avec gestion d'un état \\
        OpenAI Agents SDK & Agents, transferts entre agents et traçage \\
        Microsoft Agent Framework & Workflows d'entreprise et orchestration fan-out/fan-in \\
        Strands Agents & Composition légère d'agents et graphes d'agents \\
        PydanticAI & Agents typés et sorties structurées \\
        AutoGen & Conversations et équipes d'agents \\
        Semantic Kernel & Agents, services et plugins d'entreprise \\
        Haystack & Pipelines de composants et workflows d'agents \\
        \bottomrule
    \end{tabularx}
\end{table}

Afin d'assurer la traçabilité et la reproductibilité de la sélection,
nous relevons, pour chaque framework, la version étudiée, l'URL
officielle, la date de publication de la version, le paradigme
représenté et la raison de son inclusion. Le
Tableau~\ref{tab:metadonnees-frameworks} présente la structure utilisée
pour collecter ces informations.

\begin{table*}[t]
    \centering
    \caption{Métadonnées documentant la sélection des frameworks.}
    \label{tab:metadonnees-frameworks}
    \small
    \begin{tabular}{
        @{}
        p{0.18\textwidth}
        p{0.09\textwidth}
        p{0.10\textwidth}
        p{0.12\textwidth}
        p{0.36\textwidth}
        @{}
    }
        \hline
        \textbf{Framework}
        & \textbf{Version}
        & \textbf{Date}
        & \textbf{Source}
        & \textbf{Raison de l'inclusion} \\
        \hline

        CrewAI
        & \emph{À renseigner}
        & \emph{À renseigner}
        & \emph{À renseigner}
        & Orchestration fondée sur les rôles et les tâches \\

        Google ADK
        & \emph{À renseigner}
        & \emph{À renseigner}
        & \emph{À renseigner}
        & Agents hiérarchiques et agents de workflow \\

        LangGraph
        & \emph{À renseigner}
        & \emph{À renseigner}
        & \emph{À renseigner}
        & Workflows fondés sur des graphes avec état \\

        OpenAI Agents SDK
        & \emph{À renseigner}
        & \emph{À renseigner}
        & \emph{À renseigner}
        & Transferts entre agents et traçage natif \\

        Microsoft Agent Framework
        & \emph{À renseigner}
        & \emph{À renseigner}
        & \emph{À renseigner}
        & Orchestration de workflows d'entreprise \\

        Strands Agents
        & \emph{À renseigner}
        & \emph{À renseigner}
        & \emph{À renseigner}
        & Composition légère fondée sur des graphes \\

        PydanticAI
        & \emph{À renseigner}
        & \emph{À renseigner}
        & \emph{À renseigner}
        & Agents typés et sorties structurées \\

        AutoGen
        & \emph{À renseigner}
        & \emph{À renseigner}
        & \emph{À renseigner}
        & Équipes conversationnelles d'agents \\

        Semantic Kernel
        & \emph{À renseigner}
        & \emph{À renseigner}
        & \emph{À renseigner}
        & Architecture orientée services et plugins \\

        Haystack
        & \emph{À renseigner}
        & \emph{À renseigner}
        & \emph{À renseigner}
        & Composition d'agents orientée pipeline \\

        \hline
    \end{tabular}
\end{table*}


\section{Évaluation de la couverture fonctionnelle}
\label{sec:couverture-fonctionnelle}

La couverture fonctionnelle de GEAR est examinée dans deux directions
complémentaires :

\begin{itemize}
    \item \textbf{de GEAR vers les frameworks :} parmi les
    fonctionnalités spécifiées dans GEAR, lesquelles peuvent être
    correctement générées et exécutées dans chaque framework cible ?

    \item \textbf{des frameworks vers GEAR :} parmi les fonctionnalités
    pertinentes proposées par chaque framework, lesquelles peuvent
    être représentées par le métamodèle de GEAR ?
\end{itemize}

La première direction évalue les capacités de génération et de
préservation sémantique des connecteurs de GEAR. La seconde évalue
l'expressivité du métamodèle de GEAR. L'étude de ces deux directions
évite que GEAR obtienne artificiellement une couverture élevée en
définissant uniquement un ensemble restreint de fonctionnalités.


\subsection{Étape 1 : définition du périmètre fonctionnel}
\label{subsec:perimetre-fonctionnel}

L'évaluation se concentre sur les capacités à la conception et à
l'exécution qui sont directement liées à la spécification, à la
configuration, à l'orchestration, à la génération et à l'exécution des
\gls{mas} basés sur des \gls{llm}.

Les services de déploiement propres à un fournisseur, les
fonctionnalités cloud propriétaires, les interfaces graphiques et les
capacités de traitement de données qui ne relèvent pas directement du
cœur de GEAR sont exclues du périmètre.

Le Tableau~\ref{tab:perimetre-fonctionnel} présente les catégories
fonctionnelles étudiées.

\begin{table}[t]
    \centering
    \caption{Catégories fonctionnelles considérées dans l'évaluation de la couverture.}
    \label{tab:perimetre-fonctionnel}
    \small
    \begin{tabularx}{\columnwidth}{@{}p{0.25\columnwidth}>{\raggedright\arraybackslash}X@{}}
        \toprule
        \textbf{Catégorie} & \textbf{Exemples de fonctionnalités} \\
        \midrule
        Agent & Rôle, instructions, modèle, fournisseur et délégation \\
        Tâche & Description, entrée, sortie attendue et dépendances \\
        Outils & Déclaration, association aux agents et appels d'outils \\
        Mémoire & Mémoire individuelle et mémoire partagée \\
        Sorties & Texte, JSON et sorties structurées \\
        Orchestration & Exécution séquentielle, exécution parallèle et fan-out/fan-in \\
        Contrôle & Boucles, conditions, branchements et routage \\
        Workflow & Nœuds, arêtes, dépendances et détection des cycles \\
        Exécution & Asynchronisme, gestion des erreurs et limites d'exécution \\
        Observabilité & Traces, builds, exécutions, journaux et métriques \\
        \bottomrule
    \end{tabularx}
\end{table}

Chaque fonctionnalité est associée à une définition explicite afin de
limiter les ambiguïtés au moment de la classification. Par exemple,
nous définissons l'exécution parallèle de la manière suivante :

\begin{quote}
L'exécution parallèle désigne la capacité à exécuter simultanément au
moins deux agents ou deux tâches sans imposer de relation de dépendance
séquentielle entre eux.
\end{quote}


\subsection{Étape 2 : construire un catalogue indépendant des fonctionnalités}
\label{subsec:catalogue-fonctionnalites}

L'objectif de cette étape est de déterminer précisément les
fonctionnalités qui seront évaluées. Le catalogue ne doit pas être
construit uniquement à partir des fonctionnalités déjà proposées par
GEAR. Une telle approche pourrait conduire à une évaluation circulaire :
GEAR serait évalué uniquement par rapport à ce qu'il sait déjà faire.

Nous construisons donc un catalogue commun à partir des sources
suivantes :

\begin{itemize}
    \item le métamodèle, les schémas, les modèles de fonctionnalités
    et la documentation de GEAR ;
    \item la documentation officielle des dix frameworks ;
    \item les API publiques des frameworks ;
    \item les exemples officiellement fournis par les frameworks.
\end{itemize}

La construction du catalogue comprend quatre activités principales.


\subsubsection{Identifier les fonctionnalités de GEAR}

Nous recensons d'abord les fonctionnalités définies dans GEAR, telles
que :

\begin{itemize}
    \item la définition des agents et de leurs instructions ;
    \item la définition des tâches et de leurs dépendances ;
    \item la configuration des modèles et des fournisseurs ;
    \item l'association d'outils et de mémoires aux agents ;
    \item la composition de workflows séquentiels ;
    \item l'exécution parallèle et la synchronisation fan-out/fan-in ;
    \item la définition de boucles bornées ;
    \item la génération, l'exécution et le traçage des systèmes
    générés.
\end{itemize}


\subsubsection{Identifier les fonctionnalités des frameworks}

Nous examinons ensuite chaque framework afin d'identifier les
fonctionnalités qui relèvent de la conception et de l'orchestration des
systèmes multi-agents.

Par exemple, l'analyse de LangGraph peut faire apparaître les
fonctionnalités suivantes :

\begin{itemize}
    \item workflows représentés sous forme de graphes ;
    \item état partagé entre les nœuds ;
    \item arêtes conditionnelles ;
    \item boucles ;
    \item persistance de l'état ;
    \item interruptions humaines.
\end{itemize}

De la même manière, l'analyse de CrewAI peut faire apparaître :

\begin{itemize}
    \item agents définis par des rôles ;
    \item tâches et dépendances entre tâches ;
    \item équipes d'agents ;
    \item processus séquentiels ou hiérarchiques ;
    \item délégation ;
    \item outils et mémoire.
\end{itemize}


\subsubsection{Normaliser les fonctionnalités}

Les frameworks peuvent utiliser des termes différents pour désigner une
même capacité. Nous regroupons donc les termes équivalents sous une
fonctionnalité commune et indépendante des frameworks.

Le Tableau~\ref{tab:normalisation-fonctionnalites} présente un exemple
de cette normalisation.

\begin{table*}[t]
    \centering
    \caption{Exemple de normalisation des fonctionnalités provenant de
    différents frameworks.}
    \label{tab:normalisation-fonctionnalites}
    \begin{tabularx}{\textwidth}{@{}l l >{\raggedright\arraybackslash}X@{}}
        \toprule
        \textbf{Framework}
        & \textbf{Terme utilisé}
        & \textbf{Fonctionnalité normalisée} \\
        \midrule
        Google ADK
            & \texttt{ParallelAgent}
            & Exécution parallèle \\
        LangGraph
            & Branches parallèles
            & Exécution parallèle \\
        CrewAI
            & Tâches asynchrones
            & Exécution parallèle \\
        Strands Agents
            & Exécution concurrente
            & Exécution parallèle \\
        \bottomrule
    \end{tabularx}
\end{table*}

Chaque fonctionnalité normalisée est ensuite associée à une définition
précise. Par exemple :

\begin{quote}
L'exécution parallèle désigne la capacité à exécuter simultanément au
moins deux agents ou deux tâches indépendants, sans imposer entre eux
une relation d'ordre séquentielle.
\end{quote}


\subsubsection{Limiter le catalogue au périmètre de l'étude}

Toutes les fonctionnalités identifiées dans les documentations ne sont
pas nécessairement retenues. Nous excluons notamment :

\begin{itemize}
    \item les services de déploiement propres à un fournisseur ;
    \item les fonctionnalités cloud propriétaires ;
    \item les mécanismes de facturation ou d'authentification ;
    \item les interfaces graphiques propres aux frameworks ;
    \item les capacités de traitement de données sans lien direct avec
    la conception ou l'orchestration des agents.
\end{itemize}

Soit $\mathcal{F}_{G}$ l'ensemble des fonctionnalités de GEAR et
$\mathcal{F}_{j}$ l'ensemble des fonctionnalités identifiées pour le
framework $j$. Le catalogue initial est défini par :

\begin{equation}
\mathcal{F}
=
\mathcal{F}_{G}
\cup
\bigcup_{j=1}^{m}\mathcal{F}_{j},
\label{eq:union-fonctionnalites}
\end{equation}

où $m=10$ correspond au nombre de frameworks étudiés.

Seules les fonctionnalités appartenant au périmètre fonctionnel
$\mathcal{S}$ sont conservées dans le catalogue final :

\begin{equation}
\mathcal{F}^{*}
=
\left(
\mathcal{F}_{G}
\cup
\bigcup_{j=1}^{m}\mathcal{F}_{j}
\right)
\cap \mathcal{S}.
\label{eq:catalogue-final}
\end{equation}

Le résultat de cette étape est un catalogue structuré comportant, pour
chaque fonctionnalité, un identifiant, une catégorie, une définition
et les technologies dans lesquelles elle apparaît.

\begin{table*}[t]
    \centering
    \caption{Exemple de structure du catalogue fonctionnel.}
    \label{tab:catalogue-fonctionnel}
    \begin{tabularx}{\textwidth}{@{}l l l >{\raggedright\arraybackslash}X@{}}
        \toprule
        \textbf{ID}
        & \textbf{Catégorie}
        & \textbf{Fonctionnalité}
        & \textbf{Définition} \\
        \midrule
        AG1
            & Agent
            & Instructions
            & Définition du comportement attendu d'un agent \\
        TK1
            & Tâche
            & Dépendance
            & Utilisation du résultat d'une tâche par une autre tâche \\
        OR1
            & Orchestration
            & Séquence
            & Exécution des agents ou des tâches selon un ordre défini \\
        OR2
            & Orchestration
            & Parallélisme
            & Exécution simultanée de branches indépendantes \\
        OR3
            & Orchestration
            & Fan-in
            & Synchronisation des résultats produits par plusieurs
              branches \\
        CT1
            & Contrôle
            & Boucle bornée
            & Répétition limitée par un nombre maximal d'itérations \\
        OB1
            & Observabilité
            & Traçage
            & Enregistrement des différentes étapes d'une exécution \\
        \bottomrule
    \end{tabularx}
\end{table*}

Ce catalogue constitue le référentiel utilisé pour construire les
matrices de couverture. Il permet notamment de déterminer :

\begin{itemize}
    \item quelles fonctionnalités de GEAR peuvent être générées pour
    chaque framework ;
    \item quelles fonctionnalités des frameworks peuvent être
    représentées dans GEAR.
\end{itemize}


\subsection{Étape 3 : construction des matrices de couverture}
\label{subsec:matrices-couverture}

Deux matrices complémentaires sont construites à partir du catalogue
des fonctionnalités.

\subsubsection{Matrice de présence des fonctionnalités}

La première matrice indique si chaque fonctionnalité est proposée par
GEAR et par chacun des frameworks. Elle répond à la question suivante :
\emph{quelles fonctionnalités sont proposées par chaque technologie ?}

Le Tableau~\ref{tab:presence-fonctionnalites} illustre cette matrice.
Les points d'interrogation représentent des classifications qui devront
être établies à partir de la documentation officielle et des API.

\begin{table*}[t]
    \centering
    \caption{Exemple de matrice de présence des fonctionnalités.}
    \label{tab:presence-fonctionnalites}
    \begin{tabularx}{\textwidth}{@{}>{\raggedright\arraybackslash}Xccccc@{}}
        \toprule
        \textbf{Fonctionnalité}
        & \textbf{GEAR}
        & \textbf{CrewAI}
        & \textbf{ADK}
        & \textbf{LangGraph}
        & \textbf{\ldots} \\
        \midrule
        Exécution séquentielle
            & $\checkmark$ & $\checkmark$ & $\checkmark$
            & $\checkmark$ & \\
        Exécution parallèle
            & $\checkmark$ & $\checkmark$ & $\checkmark$
            & $\checkmark$ & \\
        Boucle bornée
            & $\checkmark$ & $\checkmark$ & $\checkmark$
            & $\checkmark$ & \\
        Routage conditionnel
            & ? & ? & ? & $\checkmark$ & \\
        Création dynamique d'agents
            & ? & ? & ? & ? & \\
        Mémoire partagée
            & $\checkmark$ & ? & ? & $\checkmark$ & \\
        \bottomrule
    \end{tabularx}
\end{table*}


\subsubsection{Matrice de préservation sémantique}

La seconde matrice indique la manière dont chaque fonctionnalité de
GEAR est implémentée dans chaque framework cible. Nous utilisons la
classification suivante :

\begin{description}
    \item[\textbf{N --- Native :}]
    GEAR utilise directement un mécanisme natif proposé par le
    framework cible.

    \item[\textbf{S --- Synthétisée :}]
    GEAR reproduit fidèlement la fonctionnalité au moyen du code
    d'orchestration généré, même si le framework cible ne propose pas
    directement une abstraction équivalente.

    \item[\textbf{A --- Adaptée :}]
    GEAR génère une implémentation, mais cette adaptation introduit une
    différence sémantique documentée.

    \item[\textbf{U --- Non supportée :}]
    la fonctionnalité ne peut pas être générée pour le framework cible.

    \item[\textbf{NA --- Non applicable :}]
    la fonctionnalité ne s'applique pas au framework considéré.
\end{description}

Le Tableau~\ref{tab:preservation-semantique} présente un exemple
illustratif.

\begin{table*}[t]
    \centering
    \caption{Exemple de matrice de préservation sémantique.}
    \label{tab:preservation-semantique}
    \begin{tabularx}{\textwidth}{@{}>{\raggedright\arraybackslash}Xcccc@{}}
        \toprule
        \textbf{Fonctionnalité GEAR}
        & \textbf{CrewAI}
        & \textbf{ADK}
        & \textbf{LangGraph}
        & \textbf{AutoGen} \\
        \midrule
        Workflow séquentiel
            & N & N & N & S \\
        Fan-out/fan-in
            & S & N & N & S \\
        Boucle bornée
            & S & N & S & A \\
        Condition d'arrêt en langage naturel
            & A & N & A & A \\
        Traçage
            & S & S & N & S \\
        \bottomrule
    \end{tabularx}
\end{table*}

Les classifications présentées dans le
Tableau~\ref{tab:preservation-semantique} sont uniquement illustratives.
La classification finale devra être établie à partir de preuves
documentaires et expérimentales.


\subsection{Étape 4 : vérifier la couverture par des preuves documentaires et expérimentales}
\label{subsec:preuves-couverture}

L'objectif de cette étape est de vérifier que les fonctionnalités
annoncées comme supportées par GEAR sont réellement représentées,
générées et exécutées conformément à leur sémantique.

Une déclaration présente dans un manifeste de connecteur, telle que :

\begin{verbatim}
parallel_workflows: true
\end{verbatim}

ne constitue pas, à elle seule, une preuve suffisante. Le code généré
peut être syntaxiquement incorrect, ne pas s'exécuter ou reproduire un
comportement différent de celui défini dans la spécification GEAR.

Chaque classification de la matrice de couverture doit donc être
étayée par cinq types de preuves.


\subsubsection{Preuve documentaire}

Nous vérifions dans la documentation officielle si le framework propose
réellement la fonctionnalité considérée.

Par exemple, la présence de \texttt{ParallelAgent} dans Google ADK
permet d'établir que ce framework propose un mécanisme natif
d'exécution parallèle.


\subsubsection{Preuve dans le connecteur GEAR}

Nous identifions le mapping ou la partie du connecteur qui transforme
la fonctionnalité GEAR en une implémentation spécifique au framework.

Par exemple, le connecteur associe le module parallèle de GEAR à
\texttt{ParallelAgent} dans Google ADK.

Dans ce cas, la fonctionnalité peut être classifiée comme
\emph{native}.

Pour un autre framework, GEAR pourrait traduire ce même module par un
appel à \texttt{asyncio.gather(...)}.

La fonctionnalité serait alors classifiée comme \emph{synthétisée}, à
condition que le comportement obtenu soit équivalent à celui attendu.


\subsubsection{Exemple minimal généré}

Pour chaque fonctionnalité, nous construisons une spécification GEAR
minimale permettant d'isoler son comportement.

Pour tester le parallélisme et la synchronisation fan-out/fan-in, nous
utilisons par exemple le workflow suivant :

\begin{equation}
A
\longrightarrow
(B \parallel C)
\longrightarrow
D.
\label{eq:workflow-parallele}
\end{equation}

Dans ce workflow :

\begin{itemize}
    \item l'agent $A$ doit être exécuté en premier ;
    \item les agents $B$ et $C$ doivent être exécutés en parallèle ;
    \item l'agent $D$ ne doit commencer qu'après la fin de $B$ et de
    $C$.
\end{itemize}

Cette spécification est générée pour chacun des frameworks étudiés.


\subsubsection{Test d'exécution}

Le code généré est exécuté avec des agents déterministes afin d'éviter
que la variabilité des \gls{llm} influence l'évaluation.

Par exemple :

\begin{itemize}
    \item l'agent $A$ enregistre le début et la fin de son exécution ;
    \item les agents $B$ et $C$ enregistrent leur démarrage, attendent
    une durée définie, puis enregistrent leur fin ;
    \item l'agent $D$ enregistre son démarrage et sa fin.
\end{itemize}

La mesure du temps peut fournir une indication sur le parallélisme,
mais elle ne constitue pas une preuve suffisante en raison des
variations possibles de l'environnement d'exécution. Nous utilisons
donc principalement les traces d'événements.


\subsubsection{Comparaison des traces}

Pour le workflow défini dans l'Équation~\ref{eq:workflow-parallele}, une
trace valide doit respecter les contraintes d'ordre suivantes :

\begin{align}
A_{\text{fin}} &< B_{\text{début}}, \\
A_{\text{fin}} &< C_{\text{début}}, \\
B_{\text{fin}} &< D_{\text{début}}, \\
C_{\text{fin}} &< D_{\text{début}}.
\end{align}

En revanche, aucun ordre particulier n'est imposé entre $B$ et $C$.
Ainsi, les deux traces suivantes sont toutes les deux valides :

\begin{verbatim}
A-debut, A-fin, B-debut, C-debut,
B-fin, C-fin, D-debut, D-fin
\end{verbatim}

\begin{verbatim}
A-debut, A-fin, C-debut, B-debut,
C-fin, B-fin, D-debut, D-fin
\end{verbatim}

Cette comparaison permet de vérifier la sémantique du workflow sans
imposer artificiellement un ordre entre les branches parallèles.


\subsubsection{Détermination de la classification}

À partir des preuves collectées, chaque fonctionnalité est classifiée
selon les règles suivantes :

\begin{description}
    \item[\textbf{Native (N) :}]
    la fonctionnalité est traduite vers un mécanisme natif équivalent
    du framework. Le code généré est valide et sa trace est conforme à
    la sémantique attendue.

    \item[\textbf{Synthétisée (S) :}]
    le framework ne propose pas directement une abstraction
    équivalente, mais le code d'orchestration généré par GEAR reproduit
    fidèlement le comportement attendu.

    \item[\textbf{Adaptée (A) :}]
    GEAR génère une implémentation utilisable, mais celle-ci introduit
    une différence sémantique documentée.

    \item[\textbf{Non supportée (U) :}]
    la fonctionnalité ne peut pas être correctement générée pour le
    framework. GEAR doit alors produire un diagnostic explicite ou
    bloquer la génération.

    \item[\textbf{Non applicable (NA) :}]
    la fonctionnalité ne concerne pas le framework considéré.
\end{description}

Par exemple, si une condition d'arrêt exprimée en langage naturel est
remplacée par un nombre fixe d'itérations, la boucle reste exécutable,
mais sa condition d'arrêt n'est pas entièrement préservée. La
fonctionnalité est donc classifiée comme \emph{adaptée}.


\subsubsection{Fiche de preuve}

Pour assurer la traçabilité de l'évaluation, chaque cellule de la
matrice de couverture est associée à une fiche de preuve. Le
Tableau~\ref{tab:fiche-preuve} présente un exemple.

\begin{table*}[t]
    \centering
    \caption{Exemple de fiche de preuve associée à une fonctionnalité.}
    \label{tab:fiche-preuve}
    \begin{tabularx}{\textwidth}{@{}l>{\raggedright\arraybackslash}X@{}}
        \toprule
        \textbf{Élément} & \textbf{Valeur} \\
        \midrule
        Fonctionnalité
            & Exécution parallèle \\
        Framework
            & Google ADK \\
        Preuve documentaire
            & Documentation officielle de \texttt{ParallelAgent} \\
        Mapping GEAR
            & Module parallèle GEAR vers \texttt{ParallelAgent} \\
        Exemple minimal
            & \texttt{parallel-fan-in.gear.yml} \\
        Test d'exécution
            & \texttt{test\_adk\_parallel\_execution} \\
        Trace attendue
            & $A$ avant $B$ et $C$ ; $B$ et $C$ avant $D$ \\
        Trace observée
            & Conforme aux contraintes d'ordre attendues \\
        Classification
            & Native (N) \\
        \bottomrule
    \end{tabularx}
\end{table*}

Le résultat de cette étape est une matrice de couverture dont chaque
classification repose sur des preuves reproductibles. L'étape 2
détermine donc \emph{ce qui doit être évalué}, tandis que l'étape 4
établit \emph{si et comment GEAR couvre effectivement chaque
fonctionnalité}.
\subsection{Étape 5 : calcul des mesures de couverture}
\label{subsec:mesures-couverture}

Soient $n_{N,j}$, $n_{S,j}$, $n_{A,j}$ et $n_{U,j}$ les nombres de
fonctionnalités respectivement classifiées comme natives, synthétisées,
adaptées et non supportées pour le framework $j$. Les fonctionnalités
classifiées comme non applicables sont exclues du dénominateur.

Soit $n_{\mathrm{app},j}$ le nombre de fonctionnalités de GEAR
applicables au framework $j$ :

\begin{equation}
n_{\mathrm{app},j}
=
n_{N,j}
+
n_{S,j}
+
n_{A,j}
+
n_{U,j}.
\label{eq:fonctionnalites-applicables}
\end{equation}


\subsubsection{Couverture de préservation sémantique}

La couverture de préservation sémantique mesure la proportion de
fonctionnalités de GEAR qui sont soit implémentées nativement, soit
fidèlement synthétisées :

\begin{equation}
C_{\mathrm{pres}}(j)
=
\frac{n_{N,j}+n_{S,j}}
     {n_{\mathrm{app},j}}.
\label{eq:couverture-preservation}
\end{equation}

Les fonctionnalités adaptées ne sont pas intégrées à cette mesure,
car leur génération introduit une différence sémantique.


\subsubsection{Couverture adaptée}

La couverture adaptée mesure la proportion de fonctionnalités qui
peuvent être générées, mais uniquement au moyen d'une adaptation
sémantique :

\begin{equation}
C_{\mathrm{adapt}}(j)
=
\frac{n_{A,j}}
     {n_{\mathrm{app},j}}.
\label{eq:couverture-adaptee}
\end{equation}


\subsubsection{Fonctionnalités non supportées}

La proportion de fonctionnalités non supportées correspond aux
fonctionnalités applicables de GEAR qui ne peuvent pas être générées
pour le framework $j$ :

\begin{equation}
C_{\mathrm{non\text{-}supp}}(j)
=
\frac{n_{U,j}}
     {n_{\mathrm{app},j}}.
\label{eq:couverture-non-supportee}
\end{equation}

Pour chaque framework, ces trois mesures vérifient :

\begin{equation}
C_{\mathrm{pres}}(j)
+
C_{\mathrm{adapt}}(j)
+
C_{\mathrm{non\text{-}supp}}(j)
= 1.
\label{eq:somme-couverture}
\end{equation}


\subsubsection{Couverture représentationnelle}

Soit $\mathcal{F}^{*}_{j}$ l'ensemble des fonctionnalités proposées
par le framework $j$ et retenues dans le périmètre de l'étude. Soit
$\mathcal{R}_{G,j}$ le sous-ensemble de ces fonctionnalités qui peuvent
être représentées par le métamodèle de GEAR.

La couverture représentationnelle de GEAR pour le framework $j$ est
définie par :

\begin{equation}
C_{\mathrm{repr}}(j)
=
\frac{
    \left|\mathcal{R}_{G,j}\right|
}{
    \left|\mathcal{F}^{*}_{j}\right|
}.
\label{eq:couverture-representationnelle}
\end{equation}

Cette mesure représente la proportion des fonctionnalités pertinentes
du framework que les utilisateurs peuvent exprimer à travers une
spécification GEAR.


\subsubsection{Couverture des diagnostics}

GEAR doit signaler explicitement les fonctionnalités adaptées et non
supportées, plutôt que de les ignorer silencieusement. Soit
$n_{\mathrm{diag},j}$ le nombre d'adaptations et d'incompatibilités
correctement détectées et signalées par GEAR.

La couverture des diagnostics est définie par :

\begin{equation}
C_{\mathrm{diag}}(j)
=
\frac{
    n_{\mathrm{diag},j}
}{
    n_{A,j}+n_{U,j}
}.
\label{eq:couverture-diagnostics}
\end{equation}

Lorsque $n_{A,j}+n_{U,j}=0$, cette mesure n'est pas applicable au
framework considéré.


\subsubsection{Couverture moyenne}

Les résultats doivent être présentés séparément pour chaque framework.
Une moyenne macro peut néanmoins être utilisée pour résumer la
couverture globale de préservation sur les $m$ frameworks :

\begin{equation}
\overline{C}_{\mathrm{pres}}
=
\frac{1}{m}
\sum_{j=1}^{m} C_{\mathrm{pres}}(j).
\label{eq:couverture-preservation-moyenne}
\end{equation}

Les résultats propres à chaque framework et la répartition entre les
fonctionnalités natives, synthétisées, adaptées et non supportées
restent cependant les résultats principaux. Nous évitons ainsi de
produire un score global pondéré qui masquerait les adaptations
sémantiques ou les fonctionnalités non supportées.


\section{Évaluation de la correction et de la fidélité des artefacts générés}
\label{sec:correction-fidelite}

L'analyse de la couverture fonctionnelle détermine quelles
fonctionnalités de GEAR peuvent être représentées dans les frameworks
cibles. Elle ne suffit cependant pas à démontrer que les artefacts
générés sont valides et qu'ils préservent effectivement le comportement
défini dans la spécification GEAR.

Nous complétons donc cette analyse par une évaluation de la correction
technique et de la fidélité sémantique des implémentations générées,
conformément à RQ2, RQ2.1 et RQ2.2.


\subsection{Corpus de workflows}
\label{subsec:corpus-workflows}

Nous construisons un corpus composé de quatre workflows représentatifs
des principales structures d'orchestration actuellement prises en
charge par GEAR :

\begin{enumerate}
    \item un workflow séquentiel ;
    \item un workflow parallèle avec synchronisation fan-out/fan-in ;
    \item un workflow comprenant une boucle bornée ;
    \item un workflow hybride combinant séquence, parallélisme et
    boucle.
\end{enumerate}

Les quatre structures utilisées dans le corpus sont résumées dans le
Tableau~\ref{tab:workflows-evaluation}.

\begin{table*}[t]
    \centering
    \caption{Workflows utilisés pour évaluer les artefacts générés.}
    \label{tab:workflows-evaluation}
    \begin{tabularx}{\textwidth}{@{}l l >{\raggedright\arraybackslash}X@{}}
        \toprule
        \textbf{ID}
        & \textbf{Structure}
        & \textbf{Propriété principalement évaluée} \\
        \midrule
        W1
            & $A \rightarrow B \rightarrow C$
            & Respect de l'ordre séquentiel et transmission des
              résultats \\
        W2
            & $A \rightarrow (B \parallel C) \rightarrow D$
            & Parallélisme, fan-out et synchronisation fan-in \\
        W3
            & $A \rightarrow [B \rightarrow C]^{k} \rightarrow D$
            & Respect du nombre maximal $k$ d'itérations \\
        W4
            & Workflow hybride
            & Combinaison de séquence, branches parallèles,
              synchronisation et boucle \\
        \bottomrule
    \end{tabularx}
\end{table*}

Chaque workflow est généré pour les dix frameworks étudiés. Le corpus
comprend ainsi :

\begin{equation}
N_{\mathrm{artefacts}}
=
N_{\mathrm{workflows}}
\times
N_{\mathrm{frameworks}}
=
4 \times 10
=
40
\end{equation}

implémentations spécifiques aux frameworks.


\subsection{Environnement d'exécution contrôlé}
\label{subsec:environnement-controle}

L'objectif de cette évaluation est de mesurer la correction du
générateur et de l'orchestration produite, et non les performances du
\gls{llm}. Nous utilisons donc des agents déterministes ou des
substituts contrôlés qui produisent toujours les mêmes résultats.

Chaque agent enregistre au minimum les événements suivants :

\begin{itemize}
    \item le début de son exécution ;
    \item la fin de son exécution ;
    \item les entrées reçues ;
    \item les résultats produits ;
    \item l'identifiant de l'itération, lorsqu'il appartient à une
    boucle.
\end{itemize}

Cette instrumentation permet de comparer directement la structure du
workflow spécifié avec le comportement de l'implémentation générée,
sans introduire la variabilité liée aux réponses des \gls{llm}.


\subsection{Niveaux de correction évalués}
\label{subsec:niveaux-correction}

Pour chaque artefact généré, nous évaluons quatre niveaux successifs de
correction.


\subsubsection{Génération}

La génération est considérée comme réussie lorsque GEAR :

\begin{itemize}
    \item valide la spécification source ;
    \item termine le processus de conversion sans erreur ;
    \item produit tous les fichiers attendus ;
    \item produit le rapport de conversion correspondant.
\end{itemize}


\subsubsection{Validité syntaxique et importation}

Un artefact est considéré comme syntaxiquement valide et importable
lorsque :

\begin{itemize}
    \item le code Python généré passe une vérification syntaxique ;
    \item les modules générés peuvent être importés ;
    \item les dépendances attendues sont correctement déclarées ;
    \item aucune référence requise n'est absente de l'artefact.
\end{itemize}


\subsubsection{Exécution}

L'exécution est considérée comme réussie lorsque :

\begin{itemize}
    \item le système généré démarre correctement ;
    \item le workflow atteint son état final ;
    \item aucune erreur non gérée n'interrompt l'exécution ;
    \item les résultats produits par les agents sont transmis aux
    étapes suivantes ;
    \item un statut de succès est enregistré dans le journal
    d'exécution.
\end{itemize}


\subsubsection{Fidélité de la trace}

La fidélité de la trace mesure si le comportement observé respecte les
contraintes définies dans la spécification GEAR. Ces contraintes
concernent notamment :

\begin{itemize}
    \item l'ordre d'exécution des agents et des tâches ;
    \item l'indépendance des branches parallèles ;
    \item la synchronisation des branches avant une étape de fan-in ;
    \item le nombre d'itérations des boucles ;
    \item la transmission des résultats entre les étapes ;
    \item l'atteinte de l'état final du workflow.
\end{itemize}


\subsection{Comparaison entre les traces attendues et observées}
\label{subsec:comparaison-traces}

Pour chaque workflow $W_i$, nous définissons un ensemble de contraintes
attendues $\mathcal{C}_{i}$. Ces contraintes sont ensuite comparées à
la trace observée $\tau_{i,j}$ pour le workflow $W_i$ généré vers le
framework $j$.

Pour le workflow parallèle :

\begin{equation}
A \rightarrow (B \parallel C) \rightarrow D,
\end{equation}

la trace observée doit respecter les contraintes suivantes :

\begin{align}
A_{\text{fin}} &< B_{\text{début}}, \\
A_{\text{fin}} &< C_{\text{début}}, \\
B_{\text{fin}} &< D_{\text{début}}, \\
C_{\text{fin}} &< D_{\text{début}}.
\end{align}

Aucune relation d'ordre n'est imposée entre $B$ et $C$. En revanche,
leurs intervalles d'exécution doivent pouvoir se chevaucher pour que
l'exécution soit considérée comme effectivement parallèle :

\begin{equation}
t_{\text{début}}(B) < t_{\text{fin}}(C)
\quad \land \quad
t_{\text{début}}(C) < t_{\text{fin}}(B).
\label{eq:chevauchement-parallele}
\end{equation}

Pour une boucle bornée comportant $k$ itérations, la trace doit
également respecter :

\begin{equation}
N_{\text{itérations observées}} = k.
\end{equation}

Enfin, chaque étape consommatrice doit recevoir les résultats produits
par les étapes dont elle dépend.


\subsection{Mesures}
\label{subsec:mesures-correction}

Soit $n$ le nombre de workflows évalués pour chaque framework. Dans
notre corpus, $n=4$.


\subsubsection{Taux de génération réussie}

Pour le framework $j$, le taux de génération réussie est défini par :

\begin{equation}
C_{\mathrm{gen}}(j)
=
\frac{
    n_{\text{générations réussies},j}
}{
    n
}.
\label{eq:taux-generation}
\end{equation}


\subsubsection{Taux de validité syntaxique et d'importation}

Le taux de validité syntaxique et d'importation est défini par :

\begin{equation}
C_{\mathrm{imp}}(j)
=
\frac{
    n_{\mathrm{artefacts\ valides\ et\ importables},j}
}{
    n
}.
\label{eq:taux-importation}
\end{equation}


\subsubsection{Taux d'exécution réussie}

Le taux d'exécution réussie est défini par :

\begin{equation}
C_{\mathrm{exec}}(j)
=
\frac{
    n_{\text{exécutions réussies},j}
}{
    n
}.
\label{eq:taux-execution}
\end{equation}


\subsubsection{Fidélité sémantique}

Soit $\mathcal{C}_{i}$ l'ensemble des contraintes attendues pour le
workflow $W_i$, et soit
$\mathcal{C}^{\mathrm{sat}}_{i,j}$ le sous-ensemble des contraintes
satisfaites par la trace observée pour le framework $j$.

La fidélité du workflow $W_i$ pour le framework $j$ est définie par :

\begin{equation}
F_{\mathrm{trace}}(i,j)
=
\frac{
    \left|\mathcal{C}^{\mathrm{sat}}_{i,j}\right|
}{
    \left|\mathcal{C}_{i}\right|
}.
\label{eq:fidelite-trace}
\end{equation}

La fidélité moyenne des traces pour le framework $j$ est ensuite
calculée par :

\begin{equation}
\overline{F}_{\mathrm{trace}}(j)
=
\frac{1}{n}
\sum_{i=1}^{n}
F_{\mathrm{trace}}(i,j).
\label{eq:fidelite-trace-moyenne}
\end{equation}

Une fidélité égale à $1$ indique que toutes les contraintes attendues
ont été respectées. Une valeur inférieure à $1$ signale qu'au moins une
propriété du workflow n'a pas été correctement préservée.


\subsubsection{Taux de correction de bout en bout}

Un artefact est considéré comme entièrement correct uniquement lorsque
les quatre conditions suivantes sont réunies :

\begin{enumerate}
    \item la génération réussit ;
    \item le code est syntaxiquement valide et importable ;
    \item l'exécution atteint son état final ;
    \item toutes les contraintes de la trace sont satisfaites.
\end{enumerate}

Le taux de correction de bout en bout pour le framework $j$ est défini
par :

\begin{equation}
C_{\mathrm{e2e}}(j)
=
\frac{
    n_{\text{artefacts entièrement corrects},j}
}{
    n
}.
\label{eq:correction-bout-en-bout}
\end{equation}


\subsection{Présentation des résultats}
\label{subsec:resultats-correction}

Les résultats sont présentés séparément pour chaque framework afin de
ne pas masquer les éventuelles différences entre les connecteurs.

\begin{table*}[t]
    \centering
    \caption{Structure du tableau de résultats pour la correction et
    la fidélité des artefacts générés.}
    \label{tab:resultats-correction}
    \begin{tabularx}{\textwidth}{@{}lccccc@{}}
        \toprule
        \textbf{Framework}
        & $\boldsymbol{C_{\mathrm{gen}}}$
        & $\boldsymbol{C_{\mathrm{imp}}}$
        & $\boldsymbol{C_{\mathrm{exec}}}$
        & $\boldsymbol{\overline{F}_{\mathrm{trace}}}$
        & $\boldsymbol{C_{\mathrm{e2e}}}$ \\
        \midrule
        CrewAI                   & -- & -- & -- & -- & -- \\
        Google ADK               & -- & -- & -- & -- & -- \\
        LangGraph                & -- & -- & -- & -- & -- \\
        OpenAI Agents SDK        & -- & -- & -- & -- & -- \\
        Microsoft Agent Framework & -- & -- & -- & -- & -- \\
        Strands Agents           & -- & -- & -- & -- & -- \\
        PydanticAI               & -- & -- & -- & -- & -- \\
        AutoGen                  & -- & -- & -- & -- & -- \\
        Semantic Kernel          & -- & -- & -- & -- & -- \\
        Haystack                 & -- & -- & -- & -- & -- \\
        \bottomrule
    \end{tabularx}
\end{table*}

Cette évaluation apporte une validation technique complémentaire à
l'analyse de la couverture fonctionnelle. La couverture indique quelles
fonctionnalités peuvent être représentées ou générées, tandis que
l'évaluation présentée dans cette section vérifie que les artefacts
produits sont effectivement valides, exécutables et conformes à la
sémantique des workflows sources.